(3 points + 3 points)
\begin{teilaufgaben}
\item
The function $u(x)$ is defined on the interval $\Omega = [-1,1].$ 

The function $u(x)$ satisfies in $\Omega$ 
\[
u''(x) - 2 \cdot u'(x) + 4 \cdot u(x)
=
4 \cdot x^2
\]
and  $u(-1) = 0, \ u(1) = 2$ on the boundary of $\Omega$. 

Determine approximate values for $u(-1/2)$, $u(0)$ and $u(1/2)$. 

Use centered finite differences with $h = 1/2$. 

\item
The function $u(x,y)$ is defined on the square $\Omega = [0, 1] \times [0,1]$.

The function $u(x,y)$ satisfies in $\Omega$
\[
\Delta u(x,y) = x^2 + y^2 - x - y
\]
and $u(x,y) = 0$ on the boundary of $\Omega$.

Determine approximate values for
$u(1/3,1/3)$, $u(2/3,1/3)$, $u(1/3,2/3)$, $u(2/3,2/3)$.

Use finite differences with $h = 1/3.$
\end{teilaufgaben}

\begin{loesung}
\begin{teilaufgaben}
\item
By the method of finite differences we have
\[
\renewcommand{\arraycolsep}{1.5pt}
\begin{array}{rclcrclcrclcr}
-4&\cdot&\tilde u(-1/2)&+&2&\cdot&\tilde u(0)&+&0&\cdot&\tilde u(1/2)&=& 1,
\\
 6&\cdot&\tilde u(-1/2)&-&4&\cdot&\tilde u(0)&+&2&\cdot&\tilde u(1/2)&=& 0,
\\
 0&\cdot&\tilde u(-1/2)&+&6&\cdot&\tilde u(0)&-&4&\cdot&\tilde u(1/2)&=&-3.
\end{array}
\]
This leads to a system of three linear equations:
\begin{equation}
\begin{pmatrix*}[r]
-4 &  2 &  0 \\
 6 & -4 &  2 \\
 0 &  6 & -4
\end{pmatrix*}
\cdot
\begin{pmatrix}
\tilde u(-1/2) \\
\tilde u(0) \\
\tilde u(1/2)
\end{pmatrix}
=
\begin{pmatrix*}[r]
 1 \\
 0 \\
-3
\end{pmatrix*}.
\tag{{\bf 2P}}
\end{equation}
Its solution is
\begin{equation}
(\tilde u(-1/2), \tilde u(0), \tilde u(1/2))
=
(-1/4 , 0, 3/4).
\tag{{\bf 1P}}
\end{equation}

\item
By the five-point-star operator we have the following system of
four linear equations: 
\begin{equation}
9
\cdot
\begin{pmatrix*}[r]
-4 &  1 &  1 &  0 \\
 1 & -4 &  0 &  1 \\
 1 &  0 & -4 &  1 \\
 0 &  1 &  1 & -4
\end{pmatrix*}
\cdot
\begin{pmatrix}
\tilde u(1/3,1/3) \\
\tilde u(2/3,1/3) \\
\tilde u(1/3,2/3) \\
\tilde u(2/3,2/3)
\end{pmatrix}
=
\begin{pmatrix}
-4/9 \\
-4/9 \\
-4/9 \\
-4/9
\end{pmatrix}.
\tag{{\bf 2P}}
\end{equation}
Hence
\begin{equation}
\tilde u(1/3,1/3)
=
\tilde u(1/3,2/3)
=
\tilde u(2/3,1/3)
=
\tilde u(2/3,2/3)
=
2/81.
\tag{{\bf 1P}}
\end{equation}
\end{teilaufgaben}
\end{loesung}

